\label{sec:identification_body}

We treat the November 30, 2022 public release of ChatGPT as a
shock to the marginal cost of obtaining private, conversational
programming assistance. We do not claim it was the only
technology changing at that date; the narrower, verifiable claim
is that it is the largest, most visible, and most
contemporaneously discussed event reducing that cost for users
whose questions overlap the range answerable by a frontier LLM.
Diffusion was fast and well-documented
\citep{eloundou2024gpts,brynjolfsson2025generative,peng2023impact,cui2024effects}.

\paragraph{What the design identifies.}
Formal pre-trend tests reject \emph{exact} parallel trends in
the $\text{AI}\times\text{Sub}$ gradient. We therefore do
\emph{not} claim a standalone natural experiment; the object we
estimate is a \emph{conditional post-ChatGPT association}
aligned with the LLM-substitutability mechanism. Its
credibility rests on a \emph{conjunction} of four diagnostics,
each failing under a different threat, so that no single
alternative survives all four: (i) the \emph{sign reversal}
between the real cutoff ($\hat\beta<0$) and 29 rolling
pre-period placebos ($\hat\beta>0$), which a smooth
continuation of the pre-trend cannot produce; (ii) the
\emph{monotone growth} of the event-study coefficient from
$-0.06$ to $-0.54$ over eight post-quarters (and on to $-0.68$
with the 2025 data), which a fading COVID-era rebound would
not generate; (iii) a \emph{pre-trend-matched} design that
restricts the comparison to tags with similar pre-treatment
slopes and returns $-0.129$ ($p=0.012$); and (iv) a HonestDiD
$\Delta^{\text{RM}}$ sensitivity bound under the full
event-study VCV, under which the estimate remains
distinguishable from zero up to a breakdown of $\bar M=0.75$.
We read $\bar M=0.75$ as a \emph{moderate} signal and lean on
the sign-reversal, monotonicity, and matched-design
evidence---none of which requires full-sample smoothness---to
carry the weight exact parallel trends cannot.

\subsection{Baseline DDD specification}\label{ssec:id_ddd_main}

The estimating equation operates at the
(tag $t$, week $w$, question type $k$) level:
\begin{equation}
\label{eq:ddd}
\begin{aligned}
\log(1+Q_{t,w,k}) =\;& \beta_{\text{DDD}}\,(\text{AI}_t \cdot \text{Post}_w \cdot \text{Sub}_k) \\
& + \gamma_1 (\text{AI}_t \cdot \text{Post}_w) + \gamma_2 (\text{AI}_t \cdot \text{Sub}_k) + \gamma_3 (\text{Post}_w \cdot \text{Sub}_k) \\
& + \alpha_{t,k} + \delta_w + \varepsilon_{t,w,k},
\end{aligned}
\end{equation}
with $\text{Post}_w = \mathbf{1}\{w \geq \text{2022-11-30}\}$,
$\text{Sub}_k \in \{0,1\}$ indicating substitutable question
types, $\alpha_{t,k}$ tag-by-question-type fixed effects, and
$\delta_w$ calendar-week fixed effects (261 weeks). Standard
errors are clustered by tag (CR1; 100 clusters)
\citep{cameron2008wild}. We estimate (\ref{eq:ddd}) under two
\emph{co-primary} treatment measures: a structural answerability
index built from four pre-treatment, non-text tag features
($-0.108$, $p=0.021$), and an embedding-based answerability
score that uses \emph{no} historical features and so is not a
relabelling of popularity ($-0.119$, $p<0.001$). We require the
result under both. Under the within-tag parallel-trends
restriction, $\beta_{\text{DDD}}$ estimates the differential
post-ChatGPT change in substitutable-type questions, relative
to non-substitutable questions in the same tags, per unit of
pre-treatment AI answerability.

\subsection{Pre-trends, placebos, and sensitivity}\label{ssec:id_event_study}\label{ssec:id_honest_did}

An event study with quarterly bins (reference bin $-1$) shows a
monotone post-period pattern, $-0.06$ at $T{+}0$ to $-0.54$ at
$T{+}8$, with pre-period coefficients negative on average but
shrinking toward zero as the cutoff approaches; formal tests
nonetheless reject exact parallel trends. Under the official
HonestDiD $\Delta^{\text{RM}}$ implementation with the full
cluster-robust VCV, the unadjusted CI is $[-0.100,-0.026]$,
remains $[-0.136,-0.010]$ at $\bar M=0.5$, and first contains
zero at $\bar M=0.75$---a moderate signal. Two placebo
exercises complement the smoothness-free reading: five
conventional pre-period cutoffs return \emph{positive}
coefficients ($+0.059$ to $+0.067$, all $p<0.012$), and a
29-step rolling-placebo sweep (mean $+0.057$) contains no
cutoff as negative as the real one. Because exact parallel
trends are rejected, we also treat a \emph{pre-trend-matched}
design as central: restricting to the 75 tags whose
pre-treatment slope lies within one IQR of the median returns
$-0.129$ ($p=0.012$), and the effect is \emph{absent} precisely
where a pre-trend artefact would be strongest (the
steepest-declining quartile gives $+0.055$, $p=0.58$), the
opposite of a mean-reversion story. Three inference variants
(two-way clustering \citep{cameron2011robust}, a wild cluster
bootstrap with 9{,}999 replicates \citep{cameron2008wild}, and
PPML on raw counts \citep{santossilva2006log}) leave the
conclusion intact. Full event-study, placebo, sensitivity, and
matched-design tables are in the replication package.

\subsection{Threats to identification}\label{ssec:id_threats}

\paragraph{COVID-driven differential mean reversion.}
If the estimate were mean reversion of a transient pandemic
boom in substitutable categories, the post-period coefficient
should fade and not accelerate. The opposite holds: the
event-study coefficient grows monotonically and steepens in
later bins, which a 2020--2021 rebound cannot produce.

\paragraph{Co-timed Stack Overflow events.}
Forces other than LLM diffusion that also post-date the cutoff
(the December 2022 AI-content ban, the mid-2023 moderator
strike, the May 2024 OpenAI licence) could in principle drive
the decline. The primary test is a donut specification dropping
the strike window (June--August 2023): the estimate is
essentially unchanged ($-0.110$, $p=0.022$). The December 2022
ban is coincident with the cutoff and absorbed by the week
fixed effects. Window-truncation estimates point \emph{against}
a one-off confound---small and imprecise early
($\hat\beta=-0.006$, $p=0.83$, ending before the strike) and
reaching magnitude later, as gradual diffusion predicts---though
this does mean onset cannot be dated to the release alone within
the original window.

\paragraph{The full-panel 2025 extension.}
The sharper discriminator is what happens \emph{after} these
one-off events. We collected calendar-2025 questions for all
100 tags from the public Stack Exchange API (registered key;
0 failed windows, 0 page-cap flags) and re-estimated on the
full 2020--2025 panel (186{,}188 cells;
Table~\ref{tab:extension_2025}), treated as a validation
extension of the 2020--2024 main panel. The pre-2025 portion
reproduces the baseline ($-0.106$); adding 2025 strengthens the
estimate to $-0.172$ ($p=0.002$), and the quarterly event-study
coefficient deepens monotonically across all of 2025, from
$-0.428$ (T$+$8) to $-0.680$ by Q4 2025 (T$+$12), with no
reversion. A transient mid-2023 episode would not produce an
association that grows for twelve post-treatment quarters;
sustained diffusion would. Exact parallel trends remain
rejected, so this is not a natural experiment, but the
full-panel evidence strongly disfavours the late-emergence and
co-timed-event readings.

\begin{table}[ht]
\centering
\caption{2025 extension on the \textbf{full 100-tag panel} (the
late-emergence test). Calendar-2025 questions for all 100 tags were
collected from the public Stack Exchange API (registered key; 0 failed
windows, 0 page-cap flags), extending the panel to 186{,}188 cells over
2020--2025. The pre-2025 portion reproduces the baseline exactly
($-0.106$, a consistency check). Adding the full 2025 year strengthens
the triple difference from $-0.108$ to $-0.172$ and tightens it to
$p=0.002$, and the quarterly event study deepens \emph{monotonically}
across all of 2025 rather than reverting---strongly disfavouring the
reading that the effect is a transient late-2024 episode (e.g.\ the
mid-2023 moderation strike) and most consistent with sustained LLM
diffusion.}
\label{tab:extension_2025}
\small
\begin{tabular}{lr}
\toprule
\multicolumn{2}{l}{\emph{Panel A. Triple difference, full 100-tag panel}} \\
\midrule
Pre-2025 portion (consistency gate)        & $-0.106$ ($p=0.023$) \\
Full 2020--2025                            & $-0.172$ ($p=0.002$) \\
Full 2020--2025, boundary-excluded         & $-0.203$ ($p=0.001$) \\
\midrule
\multicolumn{2}{l}{\emph{Panel B. Quarterly event-study coefficient, post-period}} \\
\midrule
T$+$6 (Q2 2024)  & $-0.323$ \\
T$+$8 (Q4 2024)  & $-0.428$ \\
T$+$10 (Q2 2025) & $-0.585$ \\
T$+$11 (Q3 2025) & $-0.655$ \\
T$+$12 (Q4 2025) & $-0.680$ \\
\midrule
\multicolumn{2}{l}{\emph{Panel C. Monitor persistence on new 2025 data}} \\
\midrule
Spearman(high-value drop 2024, drop 2025) & $0.903$ ($p<0.001$) \\
\bottomrule
\end{tabular}
\end{table}


\paragraph{Single non-substitutable control.}
The binary control has two categories, but only
version-environment-specific behaves as a clean placebo; the
declared boundary category advanced-architecture shows a
sizeable decline ($\approx-0.16$). Two points discipline this.
The version-environment-specific null is \emph{informative},
not merely underpowered ($-0.007$, 95\% CI $[-0.099,+0.085]$,
MDE $\approx0.13$): a substitutable-sized decline ($-0.10$ to
$-0.29$) would very likely have been detected and was not. And
the \emph{continuous} embedding substitutability score sidesteps
the dichotomy entirely, needing no clean-control category, and
still returns $-0.039$ ($p=0.004$); we treat it as the primary
safeguard. Re-assigning each category in turn as the sole
control (full table in the replication package) leaves the
estimate negative and stable under theory-led choices
(version-environment-specific as control: $-0.130$) and flips
sign only under placebo assignments that designate a
\emph{substitutable} category as the control---the
misspecification the design is built to avoid.

\paragraph{Other threats.}
The aggregate question count fell from 2.7M (2020) to 0.5M
(2024); the week fixed effect absorbs this level decline, and
$\beta_{\text{DDD}}$ is identified from \emph{differential}
movement, not levels. The pre-treatment window was not LLM-free
(Codex, Copilot preview, GPT-3 API), which biases the estimate
\emph{toward zero}, so our numbers are likely conservative with
respect to the post-2022 increment. A user-composition shift is
absorbed by tag fixed effects and the Post$\times$AI interaction
is null, so any composition shift is uniform across the tag
space rather than concentrated where answerability is highest.
